\def\topic{}
\ifdefined\MyHeaderIncluded
\else
\def\qwertyflag{true}

\ifdefined\MyHeaderIncluded

\else
  \newcommand{\foobar}{FooBar}
  \def\MyHeaderIncluded{true}
% if is till the end of the file

\documentclass[10pt,a4paper]{article}
\usepackage[utf8]{inputenc}
%\usepackage[german]{babel}
%\usepackage[cyrillic]{babel}
\usepackage[T1]{fontenc}
%\usepackage[T2A]{fontenc}
%\usepackage[english, russian]{babel}
\usepackage{amsmath}
\usepackage{amsfonts}
\usepackage{amssymb}
\usepackage{amsthm}
\usepackage{mathtools} %coloneqq

\usepackage{bussproofs} %

%\renewcommand{\thesubsection}{\arabic{subsection}}
%\newtheorem{thm}{Theorem}[section]
%\renewcommand{\thethm}{\arabic{subsection}.\arabic{thm}}

\usepackage{thmtools}
\usepackage{hyperref}
\hypersetup{
    colorlinks=true,
    linkcolor=blue,
    filecolor=magenta,      
    urlcolor=cyan,
    pdftitle={Overleaf Example},
    %pdfpagemode=FullScreen,
    }
%\usepackage{stix-type1}%\turnediota
\theoremstyle{definition}

\newtheorem{constructor}{Constructor}[section]
\newtheorem{iaxiom}{Axiom}
\newtheorem{axiom}{Axiom}
\newtheorem*{axiom*}{Axiom}
\newtheorem{definition}{Definition}[section]
\newtheorem*{definition*}{Definition}
\newtheorem{remark}{Remark}[section]
\newtheorem*{remark*}{Remark}
\newtheorem{recall}{Recall}[section]
\newtheorem*{recall*}{Recall}

\newtheorem{lemma}{Lemma}[section]
\newtheorem{theorem}{Theorem}[section]
\newtheorem*{theorem*}{Theorem}
\newtheorem{note}{Note}[section]
\usepackage{etoolbox}

%\newcommand*{\@Dashv}[2]{%
%  \reflectbox{$\m@th#1#2$}%
%}
\newcommand*{\separation}[3]{% 1)binder, 2)set, 3)condition
  \{#1\in #2: #3\}
}

% https://www.freecodecamp.org/news/simple-regex-tricks-for-beginners-3acb3fa257cb/
%  \\mathcal\{T\} |---> \topology
%  \\mathcal\{([^\}]*)\}   |--->  $1
\newcommand{\AND}{\ \ \&\ \ }
\newcommand{\any}{\mathrm{\ any}}
\newcommand{\borel}{\mathcal{B}}
\newcommand{\can}{\mathrm{can}}
\newcommand{\canonize}{\mathrm{canonize}}
\newcommand{\closure}{\mbox{cl}}
\newcommand{\complex}{\mathbb{C}}
\newcommand{\converges}{\longrightarrow}
\newcommand{\countable}{\mbox{\ countable}}
\newcommand{\ldefined}{\mbox{\ defined}}
\newcommand{\dom}{\mbox{dom}}
\newcommand{\Dom}{\mbox{Dom}}
\newcommand{\Done}{\noindent \textbf{Done.}\\}
\newcommand{\lexists}{\mbox{\ exists}}
\newcommand{\lforall}{\mbox{\ forall}}
\newcommand{\lunique}{\mbox{\ unique}}
\newcommand{\family}{\mathrm{\ family}}
\newcommand{\Fm}{\mathrm{Fm}}
\newcommand{\FuncSymb}{\mathrm{FuncSymb}}
\newcommand{\function}{\mathrm{\ function}}
\newcommand{\functional}{\mbox{\ functional}}
\newcommand{\IFF}{\ \Longleftrightarrow\ }
\renewcommand{\iff}{\ \longleftrightarrow\ }
\newcommand{\IMPLIES}{\ \Longrightarrow\ }
\renewcommand{\implies}{\ \longrightarrow\ }
\newcommand{\Inf}{\mathrm{Inf}}
\newcommand{\injective}{\mbox{\ injective}}
\newcommand{\injection}{\mbox{\ injection}}
\newcommand{\interior}{\mbox{int}}
\newcommand{\intersection}{\bigcap}
\newcommand{\isempty}{\mathrm{\ empty}}
\newcommand{\LB}{\mathrm{LB}}
\newcommand{\nghbset}{\mathcal{N}}
\newcommand{\nonempty}{\mathrm{\ nonempty}}
\newcommand{\Ord}{\mathrm{Ord}}
\newcommand{\power}{\mathcal{P}}
\newcommand{\pr}{\mathrm{pr}}
\newcommand{\fst}{\mathrm{fst}}
\newcommand{\snd}{\mathrm{snd}}
\newcommand{\PredSymb}{\mathrm{PredSymb}}
\newcommand{\Proof}{\noindent \textbf{Proof.}\\}
\newcommand{\ran}{\mathrm{ran}}
\newcommand{\Ran}{\mathrm{Ran}}
\newcommand{\real}{\mathbb{R}}
\newcommand{\relation}{\mathrm{\ relation}}
\newcommand{\set}{\mathrm{\ set}}
\newcommand{\singleton}{\mathrm{\ singleton}}
\newcommand{\semantics}{\mathrm{semantics}}
\newcommand{\semialgebra}{\mathcal{S}}
\newcommand{\Sup}{\mathrm{Sup}}
\newcommand{\surjective}{\mbox{\ surjective\ on\ }}
\newcommand{\Term}{\mathrm{Term}}
\newcommand{\Type}{\mathrm{Type}}
\newcommand{\topology}{\mathcal{T}}
\newcommand{\total}{\mbox{\ total\ on\ }}
\newcommand{\transitive}{\mathrm{\ transitive}}
\newcommand{\union}{\bigcup}
\newcommand{\lTrue}{\mathrm{True}}
\newcommand{\lFalse}{\mathrm{False}}
\newcommand{\lUncertain}{\mathrm{Uncertain}}
\newcommand{\lif}{\mathrm{if}}
\newcommand{\lthen}{\mathrm{then}}
\newcommand{\lelse}{\mathrm{else}}
\newcommand{\lexception}{\mathrm{exception}}
\newcommand{\lundef}{\mathrm{undef}}
\newcommand{\univocal}{\mathrm{\ univocal}}
\newcommand{\univocalat}{\mathrm{\ univocal\ at\ }}
\newcommand{\UB}{\mathrm{UB}}
\newcommand{\Un}{\mathcal{U}}
\newcommand{\Var}{\mathrm{Var}}

\newcommand{\notvarphi}{\phi}
\renewcommand{\phi}{\varphi}

\newcommand{\comment}[1]{\quad\mbox{(#1)}}
% local commands
%\preto\section{\undef\blipo} %etoolbox
\preto\subsection{\undef\blipo} %etoolbox
% https://tex.stackexchange.com/questions/105684/limit-macro-newcommand-scope-to-section
\usepackage{graphicx}
\usepackage{authblk}

\setlength\parindent{0pt} % whole doc of \noindent
\usepackage{soul}
%\usepackage{mathgl}

\fi % guard
\ifdefined\topic
\else
\def\topic{No Titel}
\fi
\title{\topic}
\author[1]{Georgy Dunaev}
\begin{document}
\maketitle
\fi
\section{Axioms and definitions of class theory}%Developing set-theoretic foundations}

\begin{constructor}[membership relation]
If $A$ and $B$ are classes, then $(A\in B)$ is a formula.
\end{constructor}

%\subsection{Example of a proof (MORSE-KELLEY SET THEORY, CLASSICAL SEMANTICS)}
%\begin{constructor}
%If $\varphi$ is a formula without class-quantifiers, then
%$\cc{x}{\varphi}$ is NBG-class.
%\end{constructor}

%\begin{constructor}
%If $\varphi$ is a formula, then
%$\cc{x}{\varphi}$ is MK-class.
%\end{constructor}

\begin{constructor}[subclass relation]
If $A$ and $B$ are classes, then $(A\subseteq B)$ is a formula.
\end{constructor}
\begin{definition}[subclass relation]
\[
(A\subseteq B)\iffdef(\forall x\in A:x\in B)
\]
\label{def:subclass}
\end{definition}



If the variant of logic being used doesn't have equality, then add the following constructor and definition.
\begin{constructor}[equality]%\\
If $A$ and $B$ are classes, then $(A=B)$ is a formula.
\end{constructor}
\begin{definition}[equality]
This definition is for the case if the underlying \nth{1} order logic is provided without equality.\\
If Y is not a free variable neither in A nor in B, then:
\[
(A=B) \iffdef (\forall Y: A\in Y \iff B\in Y)
\]
\end{definition}
\subsection{Axiom of extensionality}
\begin{axiom}[extensionality]
\[
(A=B)\iff ((A\subseteq B) \land (B\subseteq A))
\]
\label{ax:extensionality}
\end{axiom}
\begin{axiom*}[extensionality]
This is an other variant of the definition, with requirement that $X$ is a variable different from both $A$ and $B$). If we understand $A$ and $B$ as variables, then it is enough; if as terms, then x should not be free neither in $A$ nor in  $B$.
\[
(A=B)\iff (\forall X: X\in A \iff X\in B)
\]
\label{ax:extensionality'}
\end{axiom*}


\begin{constructor}[setness predicate]
If $A$ is a class, then $(A\set)$ is a formula.
\end{constructor}
\begin{definition}[setness predicate]
If Y is not a free variable in X, then:
\[
(X\set)\iff(\exists Y:X\in Y)
\]
\end{definition}

\subsection{Class comprehension}
\begin{constructor}[class]
If $\varphi$ is a formula, $x$ is a variable, $t$ is a term, then $\ccc{x}{t}{\varphi}$ is a class.
\end{constructor}
\begin{note}
There are no restrictions on $\varphi$ in Morse-Kelley class theory, which is being used. It is the only difference between MK and NBG: in NBG only those formulas are accepted which has only restricted quantifiers($``\forall(x\in C),\dots"$, etc).
\end{note}
%\begin{axiom}(OLD Class comprehension)
%\label{ax:comprehension}
%\[Y\in\cc{x}{\varphi} \iff (Y\set)\LAND(\varphi[Y/x])\]
%(or maybe
%\[Y\in\cc{x}{\varphi[x/Y]} \iff (Y\set)\LAND(\varphi)\]
%we need to find one expression for both cases.
%)
%\end{axiom}
\begin{definition}[closure]%\ \\
There is an order on a set of variables $\Var$. \\
Let $v_1, \dots, v_n$ be a free variables of $\phi$, a suborder.\
\
Then define the closure via the formula
\[
\text{closure}(\phi) \iffdef (\forall v_1,\dots,\forall v_n, \phi)
\]
\end{definition}

\begin{axiom}(class comprehension)%NEW 
\label{ax:comprehension}
\[
\forall(a), \forall X, (X\in\cc{X}{\varphi}) \iff ((X\set)\land \varphi)
\]
\end{axiom}

%...maybe
%\[BAD
%\cc{s}{(x\in A)} := \ccc{x}{s}{(x\in A)},
%\]
\begin{notation}
If a term is a variable, then there is no need to specify it twice:
\[
\cc{x}{\phi} \eqdef \ccc{x}{x}{\phi}.
\]
\end{notation}
\begin{notation}
If a term is value of some function $f$ (technically, it is an arbitrary term $f$), and if the argument is a variable, then it is also can be restored from the argument:
\[
\cc{f(x)}{\phi} \eqdef \ccc{x}{f(x)}{\phi}.
\]
\end{notation}

\begin{note}
%Both variants of class constructors can be expressed one from the other, since
If necessary, the class constructor can be restored from its truncated version via
\[
\ccc{x}{t}{\phi} \ceq \cc{w}{\exists x, (w=t) \land \phi}.
\]
\end{note}

\iffalse

\[
(X\in\cc{X}{\varphi}) \iff ((X\set)\land \varphi)
\]

Let's substitute X with some variable Y, which is different from X.\\
\[
(Y\in\cc{X}{\varphi}) \iff ((Y\set)\land \varphi[Y/X])
\]
Do the same step but for $\varphi[X/W]$.
\[
(X\in\cc{X}{\varphi[X/W]}) \iff ((X\set)\land \varphi[X/W])
\]
\[
(Y\in\cc{X}{\varphi[X/W]}) \iff ((Y\set)\land \varphi[X/W][Y/X])
\]

Is the following statement true:\\
"If
$\varphi[X/W]$ is defined(sic!), and if $\varphi[X/W][Y/X]$ is defined(btw the first follows from th second), and if (Y is not a free variable of $\varphi$) then $\varphi[X/W][Y/X]$ is equal to $\varphi[Y/W]$"?

...\\


Now, let's prove "old class comprehesion" statement.\\
\fi

\begin{theorem}()\\
If Y is not a free variable of $\varphi$, then:
\[
(Y\in\cc{X}{\varphi[X/W]}) \iff ((Y\set)\land(\varphi[Y/W])).
\]
\end{theorem}
\begin{proof}
($\implies$)\\
Assume $Y\in\cc{X}{\varphi[X/W]}$.\\
\dots
\end{proof}


\subsection{Definition of the empty class}

\begin{constructor}[empty class]%\ \\
The empty set symbol $\emptyclass$ is a constant closed term.
\[
(\ulcorner\emptyclass\urcorner : \Term)
\]
\end{constructor}

\begin{definition}[empty class]
\label{def:emptyclass}
\[
\emptyclass \eqdef \cc{x}{\bot}
\]
\end{definition}

\begin{definition}[empty set predicate]
\label{def:empty}
\[
(A\isempty)\iff (A=\emptyclass)
\]
\end{definition}

\begin{definition}[nonempty set predicate]
\label{def:nonempty}
\[
(A\nonempty)\iff \neg(A\isempty)
\]
\end{definition}



\subsection{Axiom of the empty set}
\begin{axiom}[the empty class is a set]
\label{ax:emptyset}
\[
(\emptyclass\set)
\]
\end{axiom}

\subsection{Definition of the universal class}
\begin{constructor}[universal class]%\ \\
The symbol $\V$ is a constant(closed term) of the language of class theory.
\end{constructor}

\begin{definition}[universal class (constant V)]
\label{def:universal class}
\[
\V\eqdef\cc{x}{\top}
\]
\end{definition}

\subsection{Definitions of unions}
\begin{definition}[elementary union]%\ \\
\label{def:union}
This is an admissible and ``elementary`` definition of a union of an arbitrary class A, which is proved to be a theorem later in this section.%Union of an arbitrary class A.
\[
\union A \eqdef \cc{x}{\exists a\in A: x\in a}
\]
\end{definition}

\begin{definition}[union of a definable conglomerate of classes]%\ \\
\label{def:uniondefcongl}
\[
\bigcup^y_{\psi} t 
\ \eqdef\ 
\cc{x}{\exists y, \psi \land (x\in t)}
\]
Side conditions:\\
1) formula $\psi$ must not depend on variable $x$,\\
2) term $t$ must not depend on variable $x$,\\
3) variable $y$ is different from variable $x$.
%Side conditions:\\
%1) variable $x$ is not free in formula $\psi$,\\
%2) variable $x$ is not free in term $t$,\\
%3) variable $x$ is not free in variable $y$.\\
%Dependencies:\\
%1) formula $\psi$ may depend on $y$.\\
%2) term $t$ may depend on $y$.
\end{definition}

\begin{definition}[union of terms built out of the elements of a class]%\ \\
Let $y$ be a variable, let $s$ and $t$ be terms.\\
If a formula has form $``(y\in s)"$, and if $y$ is also its bound variable, then it is allowed to skip writing a variable on top.
\[
\bigcup_{y\in s} t 
\ \eqdef\ 
\bigcup^y_{y\in s} t 
\]
\end{definition}
\begin{definition}[union of the elements of a class]%\ \\
\[
\bigcup A 
\ \eqdef\ 
\bigcup_{y\in A} y 
\]
\end{definition}
\begin{theorem}[``elementary def of intersection'']%\ \\
\label{thm:union}
\[
\union A = \cc{x}{\exists a\in A: x\in a}
\]
\Proof
\[
\union A
=
\bigcup_{y\in A} y  
=
\bigcup^y_{y\in A} y 
=
\cc{x}{\exists y, (y\in A) \land (x\in y)}
=
\cc{x}{\exists a\in A: x\in a}
\]
\Done
\end{theorem}

\subsection{Definitions of intersections}
\begin{definition}[elementary intersection]%\ \\
\label{def:intersection}
This is an admissible and ``elementary`` definition of an intersection of an arbitrary class A, which is proved to be a thoerem later in this section.
\[
\intersection A \eqdef \cc{x}{\forall a\in A: x\in a}
\]
\end{definition}

\begin{definition}[intersection of a definable conglomerate of classes]%\ \\
For any formula $\psi$, any term $t$, any variables $x$ and $y$, in case all of them satisfy the side conditions, the following holds
\[
\bigcap^y_{\psi} t 
\ \eqdef\ 
\cc{x}{\forall y, \psi \implies x\in t}
\]
\underline{Side conditions:}\\
1) formula $\psi$ must not depend on variable $x$,\\
2) term $t$ must not depend on variable $x$,\\
3) variable $y$ is different from variable $x$.
\label{def:interdefcongl}
\end{definition}


\begin{notation}[intersection of terms built out of the elements of a class]%\ \\
Let $y$ be a variable, let $s$ and $t$ be terms.\\
If formula has form $``(y\in s)"$, and if $y$ is also its bound variale, then it is allowed to skip writing a variable on top.
\[
\bigcap_{y\in s} t 
\ \eqdef\ 
\bigcap^y_{y\in s} t 
\]
\end{notation}
\begin{notation}[intersection of the elements of a class]%\ \\
%\[
%\bigcap G = \cc{x}{\forall y\in G. x\in y}
%\]
\[
\bigcap A 
\ \eqdef\ 
\bigcap_{y\in A} y 
\]
\end{notation}
\begin{theorem}[``elementary def of intersection'']
\label{thm:intersection}
\[
\intersection A = \cc{x}{\forall a\in A: x\in a}
\]
\Proof
\[
\intersection A
=
\bigcap_{y\in A} y  
=
\bigcap^y_{y\in A} y 
=
\cc{x}{\forall y, (y\in A) \implies (x\in y)}
=
\cc{x}{\forall a\in A: x\in a}
\]
\Done
\end{theorem}
\subsection{Definition of powerclass}
%A theory is conservatively extended with a functional  class comprehension.
\begin{definition}[powerclass]%\ \\
For any term $X$, for any variable $x$ which is not free in $X$,
\[
\power X \eqdef \cc{x}{x\subseteq X}
\tag{def. Powerclass}
\]
\label{def:powerclass}
\end{definition}
\subsection{Axiom of heredity}
\begin{axiom}[heredity]
\label{ax:heredity}
\[
(A\set) \land (B\subseteq A) \rightarrow (B\set)
\tag{ax. Heredity}
\]
\end{axiom}
\subsection{Axiom of powerset}
\begin{axiom}[powerset]
\label{ax:powerset}
\[(A\set) \rightarrow (\power A\set)\]
\end{axiom}
\begin{note}
This is a weak version, but it is okey, since together with axiom of heredity it is equivalent to the strong version.
\end{note}

%\begin{axiom} $\in$-induction. (variant that uses formulas)
%\[\left(\forall X: (\forall Y\in X: P(Y))\implies P(X)\right)\implies (\forall X: P(X))\]
%\end{axiom}
\subsection{Axiom of regularity as epsilon-induction}
\begin{axiom}[$\in$-induction]
\label{ax:epsilon induction}
For any term $\Psi$,
%(variant that uses classes)
%\[(\forall x: x\in \power\Psi \implies x\in\Psi )\implies (V\subseteq\Psi)\]
\[
(\power\Psi\subseteq\Psi)\implies(\V\subseteq\Psi)
\tag{EpsInd}
\]
\end{axiom}


\subsection{Axiom of limitation of size}
\begin{axiom}[limitation of size (of a set)]
\[
\neg(X\set) \longleftrightarrow (\exists T: (T\text{ class-function})\land (\V\subseteq T[X]))
\tag{ax. LS}
\]
\label{ax:limitation of size}
\end{axiom}

\ifdefined\qwertyflag
\undef\qwertyflag
\end{document}
\fi